\chapter{Construction de portefeuille I}\label{ch:rs:portfolio-construction-i}

Un optimiseur moyenne--variance sans contrainte reçoit cent grandes actions, une prévision d'alpha honnête et la covariance échantillonnale de leurs six derniers mois de rendements quotidiens
rendements. Il construit un portefeuille long--short dont l'exposition brute moyenne soixante et une fois son capital, et prévoit une volatilité annuelle de 36~\% et un
ratio d'information de 4.25, et le reconstruit chaque mois pendant huit ans. Le portefeuille réalise une volatilité de 207~\% et un ratio d'information de 0.92 avant
frais~; il échange quatre-vingt-quatre fois son capital par mois et, après des frais de dix points de base par transaction, son ratio d'information vaut 0.43. Les mêmes prévisions, à travers
le modèle de risque du chapitre~24 et les contraintes qu'un portefeuille réel porte réalisent 1.77 après frais. L'optimiseur a fait ce qu'on lui demandait~: il a trouvé les directions dans lesquelles
les intrants étaient les plus erronés. Ce chapitre construit \texttt{firm.portcons}, une couche de construction posée sur \texttt{firm.portopt} du livre~4, ajoute les contraintes une à une,
et valorise chacune par son prix fictif.

\section{Moyenne--variance avec un alpha et un modèle de risque}

\begin{definition}[Optimisation moyenne--variance, paramètre d'aversion au risque, frontière efficiente]\label{def:rs:portfolio-construction-i:mv}
L'\emph{optimisation moyenne--variance}\index{optimisation moyenne--variance} choisit les poids $w$ qui maximisent $\alpha^\top w - \frac{\gamma}{2} w^\top\Sigma w$
sous contraintes, pour une prévision d'alpha $\alpha$ et une prévision de covariance $\Sigma$ sur le même horizon (Markowitz, 1952). Le
\emph{paramètre d'aversion au risque}\index{paramètre d'aversion au risque} $\gamma$ règle le compromis entre rendement espéré et variance. La \emph{frontière
efficiente}\index{frontière efficiente} est l'ensemble des portefeuilles, pour tout $\gamma$, offrant le rendement espéré le plus élevé pour leur niveau de risque.
\end{definition}

\begin{proposition}[L'optimum sans contrainte et l'optimum sous contraintes d'égalité]\label{prop:rs:portfolio-construction-i:mv}
Sans contrainte, l'optimum vaut $w^* = \Sigma^{-1}\alpha/\gamma$, avec un ratio d'information ex ante $\sqrt{\alpha^\top\Sigma^{-1}\alpha}$ quel que soit $\gamma$. Avec
des contraintes d'égalité $A^\top w = 0$, il vaut $w^* = \Sigma^{-1}(\alpha - A\lambda)/\gamma$, où $\lambda$, vecteur des prix fictifs, vérifie $A^\top\Sigma^{-1}(\alpha - A\lambda) = 0$.
\end{proposition}

\begin{proof}
En annulant le gradient de $\alpha^\top w - \frac\gamma2 w^\top\Sigma w - \lambda^\top A^\top w$, on obtient $\gamma\Sigma w = \alpha - A\lambda$~; la contrainte fixe
$\lambda$. En $w^* = \Sigma^{-1}\alpha/\gamma$, $\alpha^\top w^*/\sqrt{w^{*\top}\Sigma w^*} = \sqrt{\alpha^\top\Sigma^{-1}\alpha}$.
\end{proof}

Le portefeuille du chapitre détient les cent plus grands titres de \texttt{firm.synthmkt}, reconstruit à 95 fins de mois, de l'année~3 à l'année~10, et conservé
pendant les jours du mois suivant. Son alpha est le rendement espéré propre à chaque titre dans la simulation pour le mois (la dérive persistante et la dérive post-résultats), plus un bruit deux fois plus grand,
standardisé et mis à l'échelle selon la règle du chapitre~15, $\alpha_i = \mathrm{IC}\,\sigma_i z_i$ avec un \omterm{def:rs:anatomy-of-a-predictor:ic}{coefficient d'information} de 0.05 et $\sigma_i$ la
volatilité résiduelle du titre~; le bruit est renouvelé chaque mois. Le modèle de risque est celui du chapitre~24, à horizon mensuel~; $\gamma = 12$~; les transactions coûtent dix points de base du poids échangé.
L'alpha est volontairement la vérité augmentée d'un bruit~: le chapitre étudie ce que la construction fait à une prévision donnée, non la prévision elle-même.

\section{Le maximiseur d'erreurs}

\begin{definition}[Maximisation des erreurs]\label{def:rs:portfolio-construction-i:error}
La \emph{maximisation des erreurs}\index{maximisation des erreurs} est la tendance d'un optimiseur à charger les directions dans lesquelles ses intrants sont les plus erronés~: des actifs
dont le rendement est surestimé, la variance sous-estimée ou la corrélation sous-estimée, qui semblent les meilleures transactions parce que ce sont les pires estimations (Michaud,
1989, qui les appela les optimiseurs moyenne--variance «~maximiseurs d'erreurs d'estimation~»).
\end{definition}

L'optimiseur de l'exemple d'ouverture utilise la prévision en unités de rendement (la même échelle pour chaque titre, non divisée par le risque) et la covariance échantillonnale de 126 jours pour 100
titres. Avec à peine plus de jours que de titres, les plus petites valeurs propres de la covariance échantillonnale sont bien en deçà de leur vraie valeur, et $\Sigma^{-1}$ les amplifie~: le
portefeuille prend d'énormes positions compensatoires sur des combinaisons de titres qui ont évolué ensemble par hasard au cours des six derniers mois. Son exposition brute atteint en moyenne 61, sa
volatilité prévue de 36~\% contre 207~\% réalisée, et les deux plus grandes positions ne pèsent que 9~\% de l'exposition brute en moyenne (16~\% au plus)~: les erreurs se répartissent sur
des combinaisons couvertes plutôt que concentrées sur une paire de titres. Le ratio d'information ex ante de 4.25 est la mesure que l'optimiseur fait de ses propres erreurs. Remplacer
les intrants, plutôt que de contraindre la sortie, est la première correction~: des alphas mis à l'échelle du risque vu par l'optimiseur (chapitre~15), et un modèle de risque factoriel dont les quelques
paramètres sont estimés sur de nombreux titres (chapitre~24). Sans rien changer d'autre que la \omterm{def:rs:portfolio-construction-i:constraints}{neutralité en montant}, la volatilité prévue du portefeuille passe à 13.7~\% contre
14.0~\% réalisée, et son ratio d'information atteint 1.39 après frais.

\section{Contraintes réalistes}

\begin{definition}[Neutralité en montant, en bêta et factorielle~; limite par titre~; contrainte de liquidité~; portefeuille 130/30]\label{def:rs:portfolio-construction-i:constraints}
Un portefeuille présente une \emph{neutralité en montant}\index{neutralité en montant} lorsque la somme de ses poids longs et courts est nulle, une \emph{neutralité en bêta}\index{neutralité en bêta} lorsque son
\omterm{def:rs:performance-measurement:beta}{bêta de marché} est nul, et satisfait une \emph{contrainte de neutralité factorielle}\index{contrainte de neutralité factorielle} lorsque son exposition à un facteur du modèle de risque est nulle. Une
\emph{limite par titre}\index{limite par titre} borne le poids de chaque position~; une \emph{contrainte de liquidité}\index{contrainte de liquidité} la borne par une fraction de la
valeur quotidienne moyenne échangée du titre, rapportée au capital du portefeuille. Un \emph{portefeuille 130/30}\index{portefeuille 130/30} est long de 130~\% et court de 30~\% de son capital~:
il s'agit d'un mandat en position longue seule, auquel on accorde une marge limitée pour vendre à découvert.
\end{definition}

Les contraintes entrent dans l'ordre où un desk les rencontre. Chacune est une ligne de \texttt{rs\_portcons.build} (le \cref{lst:rs:portfolio-construction-i:build}), et
\texttt{firm.portcons} transforme l'ensemble en un seul programme quadratique pour \texttt{firm.portopt.qp}~:

\begin{center}
\small
\setlength{\tabcolsep}{4pt}%
\begin{tabular}{@{}lccccccc@{}}
\toprule
 & \multicolumn{3}{c}{ratio d'information} & & & & prévision / \\
\cmidrule(lr){2-4}
réglage cumulé & ex ante & réalisé & net & rotation & brut & TC & volatilité réalisée\\
\midrule
naïf (covariance échantillonnale) & 4.25 & 0.92 & 0.43 & 84 & 61 & 0.44 & 36~\% / 207~\%\\
neutre en montant & 1.64 & 1.89 & 1.39 & 5.80 & 4.44 & 0.97 & 13.7~\% / 14.0~\%\\
+ \omterm{def:rs:portfolio-construction-i:constraints}{neutralité en bêta}, sectorielle, de style & 1.59 & 1.72 & 1.22 & 5.70 & 4.35 & 0.92 & 13.3~\% / 13.6~\%\\
+ \omterm{def:rs:portfolio-construction-i:constraints}{limites par titre} de 3~\% & 1.41 & 1.29 & 0.83 & 2.50 & 2.21 & 0.82 & 6.4~\% / 6.6~\%\\
+ brut au plus 2 & 1.44 & 1.40 & 0.94 & 2.37 & 2.00 & 0.83 & 6.1~\% / 6.3~\%\\
+ liquidité, 5~\% de la valeur quotidienne & 1.23 & 0.87 & 0.46 & 1.55 & 1.34 & 0.71 & 4.4~\% / 4.4~\%\\
+ rotation d'au plus 0.5 par mois & 1.00 & 1.94 & 1.77 & 0.53 & 1.01 & 0.58 & 3.5~\% / 3.6~\%\\
\bottomrule
\end{tabular}
\end{center}

La rotation est unilatérale et mensuelle, exprimée en multiple du capital~; le TC est le \omterm{def:rs:from-signal-to-forecast:breadth}{coefficient de transfert} (chapitre~15)~; la limite de liquidité suppose un capital de 1~milliard de dollars. Chaque
contrainte abaisse le ratio d'information ex ante, par construction~: l'optimiseur perd de la marge. Les ratios d'information réalisés sur 7.9 ans ont des erreurs types de
environ 0.5, si bien que l'essentiel des écarts entre lignes voisines relève du bruit~; les faits robustes sont la première ligne, l'effet des frais et la dernière ligne. La neutralité
aux facteurs du modèle de risque coûte peu ex ante (de 1.64 à 1.59), car l'alpha est propre à chaque titre~; Jacobs, Levy et Starer ont montré que la \omterm{def:rs:portfolio-construction-i:constraints}{neutralité en montant} et
en bêta sont en général sous-optimales, et ici elles sont presque gratuites. Les \omterm{def:rs:portfolio-construction-i:constraints}{limites par titre} et la limite d'exposition brute divisent par deux l'échelle du portefeuille et ses coûts. La limite de liquidité est saturée pour 28 des cent titres en moyenne et abaisse le ratio d'information ex ante de 1.44 à 1.23. Jagannathan et Ma ont expliqué pourquoi des contraintes fausses sur le papier aident en
pratique~: elles agissent comme un rétrécissement de l'erreur d'estimation, le carburant du maximiseur d'erreurs.

\begin{omfigure}
\begin{tikzpicture}
\begin{axis}[width=13.5cm,height=5.6cm,ybar,bar width=6pt,ylabel={ratio d'information},xtick={0,1,2,3,4,5,6},
  xticklabels={naïf,{neutre en\\montant},facteurs,titres,brut,liquidité,rotation},xticklabel style={align=center},tick label style={font=\small},label style={font=\small},
  xmin=-0.6,xmax=6.6,ymin=0,ymax=4.5,grid=major,grid style={gray!25},legend columns=3,
  legend style={font=\footnotesize,at={(0.5,-0.3)},anchor=north,draw=none,fill=none,/tikz/every even column/.append style={column sep=8pt}},
  table/col sep=comma]
\addplot[fill=gray!40,draw=gray] table[x=level,y=exante]{figdata/research/25-portfolio-construction-i/ir.csv};
\addplot[fill=omDef!40,draw=omDef] table[x=level,y=realised]{figdata/research/25-portfolio-construction-i/ir.csv};
\addplot[fill=omDef!85,draw=omDef] table[x=level,y=net]{figdata/research/25-portfolio-construction-i/ir.csv};
\legend{ex ante,réalisé,réalisé après frais}
\end{axis}
\end{tikzpicture}

\omcaption{Ratios d'information du portefeuille de cent titres à mesure que les contraintes sont ajoutées (chaque étiquette nomme la contrainte ajoutée à celles situées à sa gauche), 95 rééquilibrages mensuels. Les ratios ex ante baissent à chaque contrainte~; les ratios réalisés
ne suivent pas, et la limite de rotation donne le meilleur portefeuille après frais. Données~: \texttt{rs\_portcons.summary}.}\label{fig:rs:portfolio-construction-i:ir}
\end{omfigure}

\section{Contrôle de la rotation}

\begin{definition}[Pénalité de rotation]\label{def:rs:portfolio-construction-i:turnover}
Une \emph{pénalité de rotation}\index{pénalité de rotation} (turnover penalty) retranche $\kappa\,\lVert w - w_0\rVert_1$ de la fonction objectif, un coût linéaire sur le poids échangé à partir du
portefeuille courant $w_0$~; une limite de rotation borne plutôt $\lVert w - w_0\rVert_1$.
\end{definition}

La limite de rotation est la contrainte qui rapporte. Elle ramène les échanges mensuels du portefeuille de 1.55 à 0.53 fois le capital et ses coûts de 1.9~\% à 0.6~\% par an~; son
ratio d'information réalisé passe de 0.87 à 1.94 avant frais, ce que les frais ne sauraient expliquer. La raison tient à la prévision~: la dérive injectée persiste pendant des années,
alors que le bruit qui s'y ajoute est nouveau chaque mois. Un portefeuille qui ne peut se rapprocher que partiellement de la cible de chaque mois moyenne plusieurs prévisions mensuelles, et
cette moyenne élimine le bruit. Le \omterm{def:rs:from-signal-to-forecast:breadth}{coefficient de transfert}, la corrélation entre les poids ajustés du risque du portefeuille et la prévision (Clarke, de Silva et Thorley), passe
de 0.97 à 0.58, et cette baisse est justifiée~: la prévision qu'il transfère est plus bruitée que le portefeuille. Une pénalité ou une limite de rotation est donc deux choses à la fois, un modèle de coût
et un lisseur de prévisions~; la règle multipériode du chapitre~26 rend explicite la seconde, et le chapitre~27 évalue correctement la première.

\begin{omfigure}
\begin{tikzpicture}
\begin{axis}[width=11cm,height=5.8cm,xlabel={années depuis le premier rééquilibrage},ylabel={rendement cumulé à 10~\% de volatilité},tick label style={font=\small},
  label style={font=\small},xmin=0,xmax=8,ymin=-0.1,ymax=1.45,grid=major,grid style={gray!25},legend pos=north west,legend style={font=\footnotesize},
  yticklabel style={/pgf/number format/fixed,/pgf/number format/precision=1},scaled y ticks=false,table/col sep=comma]
\addplot[black!50,thick,dashed] table[x=year,y=naive]{figdata/research/25-portfolio-construction-i/cum.csv};
\addplot[omDef!60,thick] table[x=year,y=dollar]{figdata/research/25-portfolio-construction-i/cum.csv};
\addplot[black!70,thick] table[x=year,y=gross]{figdata/research/25-portfolio-construction-i/cum.csv};
\addplot[omDef,very thick] table[x=year,y=turnover]{figdata/research/25-portfolio-construction-i/cum.csv};
\legend{naïf,neutre en montant,+ limite brute,+ limite de rotation}
\end{axis}
\end{tikzpicture}

\omcaption{Rendements cumulés après frais, chaque portefeuille étant ramené à une volatilité annuelle de 10~\% pour que les pentes comparent les ratios d'information. Données~:
\texttt{rs\_portcons.run}.}\label{fig:rs:portfolio-construction-i:cum}
\end{omfigure}

\section{Prix fictifs et coût de chaque contrainte}

Toute contrainte a un prix fictif, le taux auquel la fonction objectif s'améliorerait si on la relâchait (livre~4, chapitre~23). \texttt{firm.portcons} le renvoie, et
davantage~: le vecteur par lequel chaque contrainte infléchit l'optimum.

\begin{proposition}[Décomposition du ratio d'information par contrainte]\label{prop:rs:portfolio-construction-i:decomp}
À l'optimum d'un problème moyenne--variance sous contraintes, $\alpha - \gamma\Sigma w^* = \sum_k g_k$, où $g_k$ rassemble les lignes de la contrainte $k$ pondérées par leurs
multiplicateurs. Le carré du ratio d'information ex ante sans contrainte, $\mathrm{IR}^2 = \alpha^\top\Sigma^{-1}\alpha$, dépasse $\gamma^2 w^{*\top}\Sigma w^*$ exactement de $\sum_k \delta_k$, avec
$\delta_k = g_k^\top\Sigma^{-1}(2\alpha - G)$ et $G = \sum_k g_k$.
\end{proposition}

\begin{proof}
La première affirmation traduit la stationnarité du système de Karush--Kuhn--Tucker restreint aux poids (les variables scindées pour les valeurs absolues contribuent
par les lignes qui les relient à $w$). Ensuite $\gamma^2 w^{*\top}\Sigma w^* = (\alpha - G)^\top\Sigma^{-1}(\alpha - G) = \alpha^\top\Sigma^{-1}\alpha - \sum_k
g_k^\top\Sigma^{-1}(2\alpha - G)$.
\end{proof}

La quantité $\gamma^2 w^\top\Sigma w$ est le carré du ratio d'information de l'alpha que les contraintes laissent à l'optimiseur, et $\delta_k$ est la part
de l'information de la prévision que la contrainte $k$ retire~; les parts s'additionnent (voir le \cref{lst:rs:portfolio-construction-i:decomp}). Pour le portefeuille entièrement contraint,
en moyenne sur les 95 mois, la limite de rotation retire 44.0~\% de l'$\mathrm{IR}^2$ sans contrainte, la liquidité 39.4~\%, les neutralités 7.0~\%, les \omterm{def:rs:portfolio-construction-i:constraints}{limites par titre} 3.1~\%~;
la limite brute et la \omterm{def:rs:portfolio-construction-i:constraints}{neutralité en montant} ne retirent rien, car une fois que la liquidité et la rotation ont réduit le portefeuille, son exposition brute vaut 1.0 et ne devient jamais contraignante. Un mois moyen,
28 bornes de liquidité sont saturées, 2.5 \omterm{def:rs:portfolio-construction-i:constraints}{limites par titre}, et la limite de rotation toujours. Ex ante, la rotation et la liquidité coûtent le plus~; ex post, la limite de rotation a rapporté plus qu'elle n'a coûté
(voir la \cref{fig:rs:portfolio-construction-i:decomp}). C'est le résultat annoncé, et la raison de présenter les deux~: les prix fictifs disent ce qu'une contrainte coûte si la
prévision est juste, et un \omterm{def:rs:vectorised-backtests:backtest}{backtest} après frais dit si elle l'était.

\begin{omfigure}
\begin{tikzpicture}
\begin{axis}[xbar,width=9.5cm,height=4.8cm,bar width=9pt,xlabel={part de l'$\mathrm{IR}^2$ ex ante sans contrainte},symbolic y coords={factor neutral,
  name limits,gross limit,liquidity,turnover limit},ytick=data,yticklabels={neutralité factorielle,limites par titre,limite brute,liquidité,limite de rotation},tick label style={font=\small},label style={font=\small},xmin=0,xmax=0.5,
  grid=major,grid style={gray!25},xticklabel style={/pgf/number format/fixed,/pgf/number format/precision=1},scaled x ticks=false,table/col sep=comma]
\addplot[fill=omDef!55,draw=omDef] table[y=group,x=share]{figdata/research/25-portfolio-construction-i/decomp.csv};
\end{axis}
\end{tikzpicture}

\omcaption{Décomposition par prix fictifs du portefeuille entièrement contraint~: la part moyenne du carré du ratio d'information ex ante sans contrainte que chaque contrainte
retire. Données~: \texttt{rs\_portcons.decomposition}.}\label{fig:rs:portfolio-construction-i:decomp}
\end{omfigure}

\section{Résolution à grande échelle}

\texttt{firm.portcons} écrit le problème sous la forme factorielle du modèle de risque, $\Sigma = XFX^\top + D$, et n'ajoute de variables scindées que pour les valeurs absolues
dont ont besoin les termes d'exposition brute et de rotation, si bien qu'une centaine de titres se ramène à cinq cents variables au plus. \texttt{firm.portopt.qp} est un solveur dense à points intérieurs, dont le coût
croît avec le cube du nombre de variables~: suffisant pour une centaine de titres et 95 mois en moins d'une minute, mais pas pour le portefeuille d'une firme de trois mille titres reconstruit
chaque jour, ce qui exige d'exploiter la structure. En introduisant les \omterm{def:rs:risk-models:families}{expositions aux facteurs}
$y = X^\top w$ comme variables fait du terme quadratique $y^\top F y + w^\top D w$ une matrice diagonale plus un petit bloc dense, et un solveur creux ou une méthode de découpage d'opérateurs
(l'ADMM du livre~4) passe alors à l'échelle avec le nombre de titres plutôt qu'avec son cube. Les démarrages à chaud à partir de la solution du mois précédent, et le filtrage des titres dont
l'alpha ne peut couvrir leur coût, font le reste.

\section{Tutoriel~: le portefeuille soixante-et-une-fois}\label{tut:rs:portfolio-construction-i:tut}

\begin{tutorial}
\textbf{Objectif.} Construire le portefeuille de cent titres avec sept ensembles cumulés de contraintes, 95 mois, et enregistrer les ratios d'information ex ante et réalisés, la rotation,
\omterm{def:rs:from-signal-to-forecast:breadth}{coefficient de transfert}, prix fictifs et la décomposition. \textbf{État final~:} le tableau, les \cref{fig:rs:portfolio-construction-i:ir,fig:rs:portfolio-construction-i:cum,fig:rs:portfolio-construction-i:decomp}.
\begin{tutsteps}
\item \textbf{Le problème}, un ensemble de contraintes à la fois.
\omcode{code/research/25-portfolio-construction-i/python/rs_portcons.py}{79}{106}{Les ensembles de contraintes du chapitre.}\label{lst:rs:portfolio-construction-i:build}
\item \textbf{Vecteurs de prix fictifs et décomposition.}
\omcode{code/firm/portcons/firm_portcons.py}{147}{170}{L'effet de chaque contrainte sur l'optimum, et sa part du ratio d'information.}\label{lst:rs:portfolio-construction-i:decomp}
\item \textbf{Exécuter} \texttt{rs\_portcons.summary(k)} pour $k = 0, \dots, 6$, \texttt{decomposition(6)} et \texttt{fig\_portcons.py}.
\end{tutsteps}
\textbf{Pour aller plus loin.} Remplacer la limite de rotation par une pénalité et tracer le ratio d'information net en fonction de $\kappa$~; rendre le bruit persistant (le même
tirage chaque mois) et constater que la limite de rotation perd son avantage.
\end{tutorial}

\section{Construction~: la couche de construction}\label{bld:rs:portfolio-construction-i:portcons}

\begin{build}
\bfield{Rôle} Chaque portefeuille que la firme négocie est construit par une seule couche~: l'alpha, le modèle de risque et les coûts en entrée, les contraintes nommées, les transactions et le prix de
chaque contrainte en sortie.
\bfield{Interface} \texttt{Problem(alpha, X, F, spec, gamma, w0)}~; \texttt{.equality(row, value, name)}, \texttt{.neutral\_factors(cols, names)},
\texttt{.bounds(lo, hi, name)}, \texttt{.liquidity(adv, participation, capital, name)}, \texttt{.gross(limit, name)}, \texttt{.turnover(limit, name)},
\texttt{.turnover\_penalty(kappa)}~; \texttt{.solve()} renvoie les poids, le ratio d'information ex ante, les prix fictifs, les nombres de contraintes saturées et les vecteurs d'effet~;
\texttt{.ir\_decomposition(result)}~; \texttt{ir\_ex\_ante}~; \texttt{trade\_list(w, w0, capital, prices, names)}.
\bfield{Règles} Les alphas sont normalisés par le risque que voit l'optimiseur~; la covariance est un modèle factoriel, jamais une covariance empirique portant sur plus de titres qu'elle ne compte de périodes~;
chaque contrainte est nommée, et son prix fictif et sa part sont présentés avec le portefeuille.
\bfield{Tests d'acceptation} \texttt{code/firm/portcons/tests/}~: la forme fermée sans contrainte~; neutralités respectées, et multiplicateurs égaux à ceux de KKT~; bornes,
limites brute et de rotation respectées et saturées~; le prix fictif de la limite brute égal à la différence finie de l'objectif~; une grande \omterm{def:rs:portfolio-construction-i:turnover}{pénalité de rotation} maintient le portefeuille immobile~;
bornes de liquidité~; vecteurs d'effet dont la somme vaut $\alpha - \gamma\Sigma w$ et une décomposition dont les parts s'additionnent~; une liste des transactions en actions.
\bfield{Pour aller plus loin} Un solveur sous forme factorielle pour des milliers de titres~; l'ajustement des facteurs propres du modèle du chapitre~24 à l'intérieur de l'optimiseur~; les lots fiscaux.
\end{build}

\begin{omsources}
\item H.~Markowitz, ``Portfolio selection'', \emph{Journal of Finance} 7(1), 1952.
\item R.~O.~Michaud, ``The Markowitz optimization enigma: is `optimized' optimal?'', \emph{Financial Analysts Journal} 45(1), 1989.
\item R.~Jagannathan and T.~Ma, ``Risk reduction in large portfolios: why imposing the wrong constraints helps'', \emph{Journal of Finance} 58(4), 2003.
\item B.~I.~Jacobs, K.~N.~Levy and D.~Starer, ``On the optimality of long--short strategies'', \emph{Financial Analysts Journal} 54(2), 1998.
\item R.~Clarke, H.~de Silva and S.~Thorley, ``Portfolio constraints and the fundamental law of active management'', \emph{Financial Analysts Journal} 58(5), 2002.
\end{omsources}

\section{Exercices}

\begin{exercise}[$\star$]\label{exo:rs:portfolio-construction-i:1}
Deux actifs non corrélés ont des alphas de 2~\% et 1~\%, et des volatilités de 20~\% et 10~\%~; $\gamma = 5$. Donner les poids non contraints et le ratio d'information ex ante.
\end{exercise}

\begin{exercise}[$\star$]\label{exo:rs:portfolio-construction-i:2}
Le carnet naïf brasse 84 fois son capital par mois à dix points de base. Quels sont ses coûts sur un an~?
\end{exercise}

\begin{exercise}[$\star$]\label{exo:rs:portfolio-construction-i:3}
Donner les erreurs types des ratios d'information réalisés de 1.77 et 1.39 mesurés sur 7.9 ans.
\end{exercise}

\begin{exercise}[$\star\star$]\label{exo:rs:portfolio-construction-i:4}
Pourquoi la covariance empirique de cent titres sur 126 jours conduit-elle à des expositions brutes de soixante fois le capital~?
\end{exercise}

\begin{exercise}[$\star\star$]\label{exo:rs:portfolio-construction-i:5}
Expliquer comment une limite de rotation peut relever le ratio d'information réalisé avant coûts.
\end{exercise}

\begin{exercise}[$\star\star$]\label{exo:rs:portfolio-construction-i:6}
Avec 1~milliard de dollars de capital et des positions limitées à 5~\% de la valeur moyenne quotidienne échangée d'un titre, quel est le poids maximal d'un titre qui s'échange pour 10~millions de dollars
par jour~?
\end{exercise}

\begin{exercise}[$\star\star\star$]\label{exo:rs:portfolio-construction-i:7}
\emph{Programmation.} Avec \texttt{rs\_portcons.decomposition(6)} et \texttt{run(6)}, expliquer pourquoi la part de la limite brute est nulle alors qu'elle figure dans le problème.
\end{exercise}

\begin{exercise}[$\star\star\star$]\label{exo:rs:portfolio-construction-i:8}
\emph{Trouver la faille.} «~L'optimiseur affiche un ratio d'information ex ante de 4.25 sur une covariance empirique des six derniers mois, nous dimensionnerons donc le fonds à partir de ce chiffre.~»
\end{exercise}

\section{Problème~: le carnet aux soixante et une fois}

\begin{problem}[{Problème du week-end --- contraintes, chiffrées}]\label{pb:rs:portfolio-construction-i:1}
Le carnet de cent titres du chapitre sur \texttt{firm.synthmkt}.

\medskip\noindent\textbf{Partie I --- L'optimiseur.}
\begin{enumerate}
\item Écrire l'optimum non contraint et son ratio d'information ex ante.
\item Comment l'alpha et le modèle de risque sont-ils construits, et pourquoi l'alpha vaut-il la vérité plus du bruit~?
\item À quoi ressemble le carnet naïf~: brut, volatilité prévue et réalisée, ratios d'information ex ante et réalisé, rotation~?
\item Pourquoi ses erreurs se répartissent-elles sur des combinaisons couvertes plutôt que sur deux titres~?
\end{enumerate}

\medskip\noindent\textbf{Partie II --- Contraintes.}
\begin{enumerate}[resume]
\item Que change le modèle de risque factoriel avec la \omterm{def:rs:portfolio-construction-i:constraints}{neutralité en montant}~?
\item Que coûtent les neutralités factorielles ex ante, et pourquoi si peu~?
\item Que font les \omterm{def:rs:portfolio-construction-i:constraints}{limites par titre} et la limite brute à l'échelle et aux coûts~?
\item Combien de bornes de liquidité sont saturées, et que coûtent-elles ex ante~?
\item Pourquoi la plupart des écarts de ratio d'information réalisé entre réglages voisins sont-ils du bruit~?
\end{enumerate}

\medskip\noindent\textbf{Partie III --- Rotation.}
\begin{enumerate}[resume]
\item Que deviennent la rotation et les coûts avec la limite de rotation~?
\item Pourquoi le ratio d'information réalisé monte-t-il avant coûts~?
\item Pourquoi le \omterm{def:rs:from-signal-to-forecast:breadth}{coefficient de transfert} baisse-t-il, et est-ce un mal~?
\end{enumerate}

\medskip\noindent\textbf{Partie IV --- Le verdict.}
\begin{enumerate}[resume]
\item Énoncer la proposition qui sous-tend la décomposition.
\item Énoncer le \emph{résultat nommé}~: la part de l'information ex ante non contrainte que prend chaque contrainte, et ce que chacune rapporte après coûts.
\item Pourquoi la limite brute et la \omterm{def:rs:portfolio-construction-i:constraints}{neutralité en montant} ne prennent-elles rien~?
\item Quel carnet faudrait-il trader, et que faudrait-il présenter avec lui~?
\item Quand une covariance empirique est-elle acceptable dans un optimiseur~?
\item Que diraient Jagannathan et Ma des \omterm{def:rs:portfolio-construction-i:constraints}{limites par titre}~?
\item Comment résoudrait-on le problème pour trois mille titres~?
\item En une phrase~: que vaut le ratio d'information ex ante d'un optimiseur~?
\end{enumerate}
\end{problem}

\section{Questions d'entretien}

\begin{interviewq}[$\star$ \iqroles{chercheur, trader}]\label{iq:rs:portfolio-construction-i:1}
Pourquoi les portefeuilles moyenne--variance non contraints donnent-ils de mauvais résultats \omterm{def:rs:overfitting:overfitting}{hors échantillon}~?
\end{interviewq}

\begin{interviewq}[$\star\star$ \iqroles{chercheur}]\label{iq:rs:portfolio-construction-i:2}
Qu'est-ce qu'un prix fictif dans la construction d'un portefeuille, et comment l'utiliseriez-vous~?
\end{interviewq}

\begin{interviewq}[$\star\star$ \iqroles{chercheur, risque}]\label{iq:rs:portfolio-construction-i:3}
Votre carnet doit être neutre en montant, en bêta et en secteur. Que coûte cette neutralité, et quand est-elle gratuite~?
\end{interviewq}

\begin{interviewq}[$\star\star$ \iqroles{trader}]\label{iq:rs:portfolio-construction-i:4}
Comment fixez-vous une limite de liquidité, et que devient le carnet quand le capital croît~?
\end{interviewq}

\begin{interviewq}[$\star\star$ \iqroles{chercheur}]\label{iq:rs:portfolio-construction-i:5}
Qu'est-ce que le \omterm{def:rs:from-signal-to-forecast:breadth}{coefficient de transfert}, et pourquoi un coefficient plus faible peut-il être meilleur~?
\end{interviewq}

\begin{interviewq}[$\star\star\star$ \iqroles{chercheur, développeur}]\label{iq:rs:portfolio-construction-i:6}
Concevez l'optimiseur d'un carnet de trois mille titres rééquilibré chaque jour~: formulation, solveur, et ce que vous précalculez.
\end{interviewq}
